% !PRINTED 8
\weberpage{8}
%<<CONT
$Q_1$ verschieden sind. Denn wäre etwa $a_1 c = a_2 b$, so würde folgen, dass $c = a_1^{-1} a_2 b$ sein müsste; es ist aber $a_1^{-1} a_2$ in $Q$ enthalten, also mit einem der Elemente $1, a_1, a_2 \ldots a_{\nu-1}$ identisch, und es wäre also $c$ gegen die Voraussetzung in $Q_1$ enthalten. So fahren wir fort, die Systeme $Q, Q_1, Q_2 \ldots Q_{\mu-1}$ zu bilden, von denen das letzte
\[
Q_{\mu-1} = g, a_1 g, a_2 g \ldots a_{\nu-1} g
\]
sei. Dann muss aber $P$ durch die Systeme $Q, Q_1, Q_2 \ldots Q_{\mu-1}$ erschöpft sein, und es folgt, da diese Systeme lauter verschiedene Elemente enthalten,
\[
n = \nu \mu,
\]
oder der Satz:

\noindent 1.\ \gesperrt{Der Grad einer Gruppe ist durch den Grad jedes ihrer Divisoren theilbar. Der Quotient $\mu = n : \nu$ soll der Index des Theilers $Q$ heissen.}

Die Systeme $Q_1, Q_2 \ldots Q_{\mu-1}$ nennen wir, wie in dem speciellen Falle der Permutationsgruppen, die zu $Q$ gehörigen \gesperrt{Nebengruppen} und bezeichnen sie durch
\[
Q_1 = Qb, \quad Q_2 = Qc \ldots Q_{\mu-1} = Qg,
\]
so dass wir auch symbolisch
\weq{1}{P = Q + Qb + Qc + \cdots + Qg}
setzen können. Bisweilen werden wir auch das ganze System $Q, Q_1 \ldots Q_{\mu-1}$ ($Q$ selbst eingeschlossen) als ein System von Nebengruppen bezeichnen, und also die Darstellung (1) die \gesperrt{Zerlegung von $P$ in ein System von Nebengruppen} nennen.

Aus der Bildungsweise der Nebengruppen folgt noch, dass, wenn $b$, $c$ irgend zwei Elemente aus $P$ sind, die beiden Systeme $Qb$ und $Qc$ entweder ganz identisch sind oder kein gemeinsames Element enthalten. Denn ist $a_1 b = a_2 c$, worin $a_1$, $a_2$ zwei Elemente aus $Q$ sind, so ist auch für jedes andere Element $a$ aus $Q$:
\[
a a_1 b = a a_2 c,
\]
und wenn $a$ die ganze Gruppe $Q$ durchläuft, so durchläuft, wie schon in \S.\,1 bemerkt, auch jedes der beiden $a a_1$ und $a a_2$ die ganze Gruppe $Q$; folglich sind $Qb$ und $Qc$ identisch.

Wie in Bd.\,I, \S.\,154, 2. können wir $P$ auch in der folgenden Art in Nebengruppen zerlegen:
\[
P = Q + b^{-1} Q + c^{-1} Q + \cdots + g^{-1} Q.
\]
